\chapter{Implémentations Algorithmiques de Référence}
\label{app:code_listings}

Cette annexe rassemble les extraits de code source fondamentaux constituant le cœur algorithmique de la plateforme CBIR.

\section{Module d'Indexation Spatiale Métrique (\texttt{cbir/indexing/vptree.py})}

\begin{lstlisting}[language=Python, caption={Implémentation intégrale de l'index VP-Tree}]
import heapq
import random
import numpy as np
from typing import List, Tuple, Optional, Callable
from dataclasses import dataclass

@dataclass
class _Node:
    point: np.ndarray
    point_id: str
    mu: float
    left: Optional['_Node'] = None
    right: Optional['_Node'] = None

class VPTree:
    """Arbre d'espace metrique Vantage-Point Tree pour recherche k-NN."""

    def __init__(self, distance_fn: Callable[[np.ndarray, np.ndarray], float]):
        self.distance_fn = distance_fn
        self.root: Optional[_Node] = None
        self.n_nodes: int = 0
        self.comparisons_last_query: int = 0

    def build(self, items: List[Tuple[str, np.ndarray]]):
        """Construit l'arbre recursivement a partir de couples (id, vecteur)."""
        self.n_nodes = len(items)
        self.root = self._build_recursive(list(items))

    def _build_recursive(self, items: List[Tuple[str, np.ndarray]]) -> Optional[_Node]:
        if not items:
            return None

        # Selection aleatoire uniforme du Vantage Point
        vp_idx = random.randint(0, len(items) - 1)
        vp_id, vp_point = items[vp_idx]
        rest = items[:vp_idx] + items[vp_idx + 1:]

        if not rest:
            return _Node(vp_point, vp_id, mu=0.0)

        # Calcul des distances au pivot
        distances = [self.distance_fn(vp_point, pt) for _, pt in rest]
        mu = float(np.median(distances))

        # Partitionnement binaire selon la mediane mu
        left_items = [item for item, d in zip(rest, distances) if d < mu]
        right_items = [item for item, d in zip(rest, distances) if d >= mu]

        node = _Node(vp_point, vp_id, mu)
        node.left = self._build_recursive(left_items)
        node.right = self._build_recursive(right_items)
        return node

    def query(self, query: np.ndarray, k: int = 10) -> List[Tuple[str, float]]:
        """Recherche des k plus proches voisins avec elagage par tas-max."""
        if self.root is None:
            return []

        # Tas-max stockant (-distance, point_id)
        heap: List[Tuple[float, str]] = []
        self.comparisons_last_query = 0

        self._query_node(self.root, query, k, heap)

        results = [(-dist, pid) for dist, pid in heap]
        results.sort(key=lambda x: x[1])
        return [(pid, dist) for dist, pid in sorted(results, key=lambda x: x[1])]

    def _query_node(self, node: Optional[_Node], query: np.ndarray, k: int, heap: List[Tuple[float, str]]):
        if node is None:
            return

        self.comparisons_last_query += 1
        dist = self.distance_fn(query, node.point)

        # Mise a jour du tas de priorite
        if len(heap) < k:
            heapq.heappush(heap, (-dist, node.point_id))
        elif dist < -heap[0][0]:
            heapq.heapreplace(heap, (-dist, node.point_id))

        tau = -heap[0][0] if len(heap) >= k else float('inf')

        # Elagage par l'inegalite triangulaire
        if dist < node.mu:
            if dist - tau < node.mu:
                self._query_node(node.left, query, k, heap)
            tau = -heap[0][0] if len(heap) >= k else float('inf')
            if dist + tau >= node.mu:
                self._query_node(node.right, query, k, heap)
        else:
            if dist + tau >= node.mu:
                self._query_node(node.right, query, k, heap)
            tau = -heap[0][0] if len(heap) >= k else float('inf')
            if dist - tau < node.mu:
                self._query_node(node.left, query, k, heap)
\end{lstlisting}

\section{Module des Métriques de Distance (\texttt{cbir/distances/metrics.py})}

\begin{lstlisting}[language=Python, caption={Bibliothèque vectorielle de calcul de distances}]
import numpy as np

def l1_distance(x: np.ndarray, y: np.ndarray) -> float:
    """Distance de Manhattan (L1)."""
    return float(np.sum(np.abs(x - y)))

def l2_distance(x: np.ndarray, y: np.ndarray) -> float:
    """Distance Euclidienne (L2)."""
    return float(np.linalg.norm(x - y))

def chi_squared_distance(x: np.ndarray, y: np.ndarray, eps: float = 1e-10) -> float:
    """Distance du Chi-Deux (Chi2) pour histogrammes."""
    num = (x - y) ** 2
    denom = x + y + eps
    return float(0.5 * np.sum(num / denom))

def bhattacharyya_distance(x: np.ndarray, y: np.ndarray, eps: float = 1e-10) -> float:
    """Distance de Bhattacharyya."""
    bc = np.sum(np.sqrt(np.clip(x * y, 0.0, None)))
    bc = np.clip(bc, eps, 1.0)
    return float(-np.log(bc))

def cosine_distance(x: np.ndarray, y: np.ndarray, eps: float = 1e-10) -> float:
    """Distance Cosinus (1 - similarite cosinus)."""
    norm_x = np.linalg.norm(x)
    norm_y = np.linalg.norm(y)
    if norm_x < eps or norm_y < eps:
        return 1.0
    similarity = np.dot(x, y) / (norm_x * norm_y + eps)
    return float(np.clip(1.0 - similarity, 0.0, 2.0))
\end{lstlisting}

\section{Module d'Extraction LBP (\texttt{cbir/descriptors/lbp.py})}

\begin{lstlisting}[language=Python, caption={Extracteur de motifs locaux binaires LBP}]
import cv2
import numpy as np
from skimage.feature import local_binary_pattern

class LBPExtractor:
    def __init__(self, n_points: int = 24, radius: int = 3, method: str = 'uniform'):
        self.n_points = n_points
        self.radius = radius
        self.method = method
        self.n_bins = n_points + 2 if method == 'uniform' else 2 ** n_points

    def extract(self, image_path: str) -> np.ndarray:
        image = cv2.imread(image_path, cv2.IMREAD_GRAYSCALE)
        if image is None:
            raise ValueError(f"Impossible de charger l'image : {image_path}")

        image = cv2.resize(image, (256, 256), interpolation=cv2.INTER_AREA)
        lbp = local_binary_pattern(image, self.n_points, self.radius, method=self.method)

        hist, _ = np.histogram(lbp.ravel(), bins=self.n_bins, range=(0, self.n_bins))
        hist = hist.astype(np.float32)
        total = hist.sum()
        if total > 0:
            hist /= total
        return hist
\end{lstlisting}
